\documentclass[11pt,a4paper]{article}
\usepackage[margin=25mm,headheight=14pt]{geometry}
\usepackage{fontspec}
\setmainfont{lmroman10-regular.otf}[BoldFont=lmroman10-bold.otf,ItalicFont=lmroman10-italic.otf,BoldItalicFont=lmroman10-bolditalic.otf]
\setsansfont{lmsans10-regular.otf}
\setmonofont{lmmono10-regular.otf}[Scale=MatchLowercase]
\newfontfamily\authorsymbol{Noto Sans Symbols}[Scale=MatchUppercase]
\usepackage{amsmath,amssymb,amsthm,mathtools,booktabs,array}
\usepackage{microtype}
\usepackage[hidelinks,unicode]{hyperref}
\usepackage{fancyhdr}
\pagestyle{fancy}\fancyhf{}
\fancyhead[L]{\small\itshape Affine fibers and a Fibonacci bound}
\fancyhead[R]{\small\thepage}\renewcommand{\headrulewidth}{0pt}
\setlength{\parindent}{1.2em}\setlength{\parskip}{0pt}
\setlength{\emergencystretch}{2em}
\newtheorem{theorem}{Theorem}
\newtheorem{proposition}[theorem]{Proposition}
\newtheorem{lemma}[theorem]{Lemma}
\newtheorem{corollary}[theorem]{Corollary}
\theoremstyle{definition}\newtheorem{definition}[theorem]{Definition}
\newcommand{\F}{\mathbb F_2}
\newcommand{\andb}{\mathbin{\&}}
\newcommand{\res}[1]{[#1]_M}
\DeclareMathOperator{\wt}{wt}
\DeclareMathOperator{\rank}{rank}
\hypersetup{pdftitle={Affine Fibers and a Fibonacci Bound for Boolean–Modular Inversion},pdfauthor={lmtlssss🜐},pdfsubject={Mathematics: discrete mathematics and exact inversion},pdfkeywords={affine fibers, binary linear maps, Fibonacci numbers, carry automata}}
\title{\LARGE Affine Fibers and a Fibonacci Bound\\for Boolean--Modular Inversion}
\author{lmtlssss{\authorsymbol 🜐}}
\date{27 October 2026}
\begin{document}
\maketitle\thispagestyle{plain}
\begin{abstract}
Let $L$ be an invertible binary-linear map on words of width $w$.
We describe every solution of $L(a+D)+L(a)=S\pmod{2^w}$ as a disjoint
union of affine subspaces over $\F$. A feasible XOR signature reduces both
modular additions to linear mask equations. Two carry automata enumerate
the input and output signatures, so the construction uses the smaller
list. Each list has at most $F_{w+1}$ entries, where $F_0=0$ and $F_1=1$.
This bound is attained for signature enumeration. Ordinary binary
elimination gives a conservative $O(w^3F_{w+1})$ bit-operation bound for
compact output. Two classes of target words require only one affine
cell. We supply the complete proof, a worked five-cell fiber, and an
executable verifier.
\end{abstract}
\noindent\textit{Mathematics Subject Classification (2020).}
05A15; 11B39, 68W30.\par
\noindent\textit{Keywords.} Affine subspaces; binary linear maps;
modular arithmetic; carry automata.\par
\noindent\textit{DOI.} \href{https://doi.org/10.5281/zenodo.23697726}{10.5281/zenodo.23697726}.

\section{The inverse problem}
Fix $w\geq1$ and put $M=2^w$. We identify $\F^w$ with the integers
$0,\ldots,M-1$, using the least significant bit as coordinate zero.
The symbols $\oplus$ and $\andb$ denote bitwise XOR and AND. An overbar
means complement within $w$ bits. Write $\res z$ for the least nonnegative
residue of $z$ modulo $M$. Addition inside an argument of $L$ is modular.
The linearity of $L$ is with respect to $\oplus$, not integer addition.

Given $D,S\in\{0,\ldots,M-1\}$ and an invertible map
$L:\F^w\to\F^w$, define
\begin{equation}\label{eq:fiber}
\mathcal A_{D,S}(L)=\{a:L(\res{a+D})+L(a)\equiv S\pmod M\}.
\end{equation}
An affine cell $b+V$ means the set $\{b\oplus v:v\in V\}$ for a
binary-linear subspace $V$. We return a base and an independent basis
for each cell. This representation specifies all its members without
listing them.

Define the two sets of feasible signatures
\begin{align}
\Delta_D&=\{a\oplus\res{a+D}:0\leq a<M\},\label{eq:deltas}\\
K_S&=\{x\oplus y:x+y\equiv S\pmod M\}.\label{eq:ks}
\end{align}
The unrestricted expression in \eqref{eq:ks} uses two words $x,y$.
Its signatures become relevant to \eqref{eq:fiber} through the relation
$k=L(\delta)$.

\begin{theorem}\label{thm:main}
For every invertible binary-linear $L$, the fiber
$\mathcal A_{D,S}(L)$ is a disjoint union of at most
\[
N_{D,S}=\min\{|\Delta_D|,|K_S|\}\leq F_{w+1}
\]
affine cells. A constructive algorithm visits exactly one signature per
entry of the selected smaller list, discards empty cells, and returns
all remaining cells by binary linear elimination. Its conservative
bit-operation cost is $O(w^3(1+N_{D,S}))$ for compact output.
\end{theorem}

Here the number of signatures is not the number of solutions. A single
cell of dimension $d$ represents $2^d$ solutions. Expanding such a cell
requires additional output work.

Finite-state carry calculations are established. Mouha et al.\ \cite{mouha}
use finite-state methods and matrix products to calculate differential
probabilities and counts of possible differences. Velichkov et al.\ \cite{velichkov}
extend that framework to compositions involving rotations. Qin and
Wu\ \cite{qin} retain dependencies between consecutive modular additions.
This note uses those familiar ingredients to state a complete affine
partition for \eqref{eq:fiber}, with a bidirectional signature enumeration
and an explicit certificate for each cell. The construction makes no
locality assumption on the rows of $L$.

\section{A signature gives linear equations}
Put $\ell=2^{w-1}-1$. For a candidate output signature $k$, define
\begin{equation}\label{eq:parameters}
\delta=L^{-1}(k),\qquad d=\res{\delta-D},\qquad e=\res{S-k}.
\end{equation}
If $d$ and $e$ are even, set $u=d/2$ and $v=e/2$. These are ordinary
integer halves of the canonical residues. Set
\[
p=\delta\andb\ell,\qquad q=\overline{k}\andb\ell.
\]

\begin{lemma}\label{lem:cell}
The solutions with signature $k$ are exactly the solutions of
\begin{equation}\label{eq:masks}
a\andb p=u,\qquad L(a)\andb q=v,
\end{equation}
provided $d,e$ are even and
\begin{equation}\label{eq:support}
u\andb\overline p=0,\qquad v\andb\overline q=0.
\end{equation}
If any of these conditions fails, the cell is empty.
\end{lemma}
\begin{proof}
The integer identity
\[
a\oplus\delta=a+\delta-2(a\andb\delta)
\]
shows that $\res{a+D}=a\oplus\delta$ if and only if
$2(a\andb\delta)\equiv\delta-D\pmod M$.
Multiplication by two removes the top bit modulo $M$. Thus this condition
is equivalent to $d$ being even and $a\andb p=u$.
The support condition on $u$ is necessary and sufficient for that mask
equation to be consistent in isolation.

Write $x=L(a)$. Under the input condition, linearity gives
$L(\res{a+D})=x\oplus k$. The second integer identity
\[
x+(x\oplus k)=k+2(x\andb\overline k)
\]
then converts \eqref{eq:fiber} into $2(x\andb\overline k)\equiv S-k\pmod M$.
This is equivalent to $e$ being even and $x\andb q=v$.
The argument proves both directions, including the discarded top carry.
\end{proof}

If $r_i$ is row $i$ of the matrix of $L$, each one-bit position of $p$
gives an equation $a_i=u_i$. Each one-bit position of $q$ gives
$\langle r_i,a\rangle=v_i$ over $\F$. These are at most $2(w-1)$ linear
equations in $w$ variables. Binary elimination returns either inconsistency
or one base with a basis of its nullspace.

\begin{proposition}\label{prop:partition}
The nonempty cells from Lemma~\ref{lem:cell} form a disjoint partition of
$\mathcal A_{D,S}(L)$ when all signatures of either $K_S$ or $\Delta_D$
are considered, with $k=L(\delta)$ in the second case.
\end{proposition}
\begin{proof}
Every solution $a$ determines the unique pair
\[
\delta=a\oplus\res{a+D},\qquad k=L(\delta).
\]
This pair belongs to $\Delta_D\times K_S$ and satisfies the two mask
equations. Conversely, Lemma~\ref{lem:cell} makes every point of a returned
cell a solution. A point in two different cells would have two different
values of $a\oplus\res{a+D}$ or of its image under $L$, which is impossible.
\end{proof}

\section{Two carry automata and the Fibonacci bound}
For the output list, let $c_i$ be the carry entering bit $i$ in $x+y=S$,
with $c_0=0$. The XOR bit is
\begin{equation}\label{eq:outbit}
k_i=S_i\oplus c_i.
\end{equation}
If $k_i=1$, the two input bits differ and $c_{i+1}=c_i$. If $k_i=0$,
the two input bits are equal and either outgoing carry is possible.
The outgoing carry at bit $w-1$ is discarded.

For the input list, let $t_i$ be the carry entering bit $i$ of $a+D$,
with $t_0=0$. Then
\begin{equation}\label{eq:inbit}
\delta_i=D_i\oplus t_i.
\end{equation}
If $\delta_i=0$, then $D_i=t_i$ and $t_{i+1}=t_i$. If $\delta_i=1$,
the bit $a_i$ selects either outgoing carry. Again the final carry is
discarded.

Each signature determines its entering carries from \eqref{eq:outbit}
or \eqref{eq:inbit}. Therefore different paths give different signatures.
Conversely, choose compatible operand bits at each branch of a path.
This produces words realizing that signature. The automata are thus
complete and duplicate-free.

\begin{lemma}\label{lem:fib}
For every $D,S$, both $|\Delta_D|$ and $|K_S|$ are at most $F_{w+1}$.
For each width, this maximum is attained separately in each automaton.
\end{lemma}
\begin{proof}
Let $(A,B)$ count partial paths with entering carry zero and one.
The initial pair is $(1,0)$. According to the prescribed bit, the pair
updates by one of the two rules
\begin{equation}\label{eq:updates}
(A,B)\longmapsto(A+B,B),\qquad
(A,B)\longmapsto(A,A+B).
\end{equation}
The assignment of a prescribed bit to a rule is reversed between the
input and output automata.

After $i$ updates, write the counts in decreasing order $R\geq T$.
Induction gives $R\leq F_{i+1}$ and $T\leq F_i$. An update has largest
entry at most $R+T\leq F_{i+2}$ and smaller entry at most
$R\leq F_{i+1}$. There are $w-1$ relevant outgoing-carry decisions.
Hence the number of signatures is at most $F_w+F_{w-1}=F_{w+1}$.
Starting with the rule that gives $(1,1)$ and alternating the two rules
attains equality. The top prescribed bit does not change the count.
\end{proof}

\begin{proof}[Proof of Theorem~\ref{thm:main}]
Count the two signature lists using \eqref{eq:updates}, then enumerate
the smaller list. Invertibility of $L$ couples $k$ and $\delta$ uniquely;
there is no Cartesian-product search over the two lists. For each
signature, apply \eqref{eq:parameters}--\eqref{eq:support} and solve the
linear system \eqref{eq:masks}. Proposition~\ref{prop:partition} proves
completeness and disjointness. Lemma~\ref{lem:fib} proves the stated
signature bound.

An inverse matrix can be prepared in $O(w^3)$ bit operations. Each
signature requires $O(w^2)$ work to apply the map or its inverse, and at
most $O(w^3)$ work for elimination. A prefix traversal generates each
signature with at most $w$ steps. The preliminary counts have $O(w)$
bits and need $O(w^2)$ bit work. These estimates give
$O(w^3(1+N_{D,S}))$ overall. The explicit bases occupy
$O(w^2N_{D,S})$ bits if all cells are retained.
\end{proof}

Since $F_{w+1}=\Theta(\varphi^w)$ for
$\varphi=(1+\sqrt5)/2$, this is an $O(w^3\varphi^w)$ compact-output
bound. The Fibonacci maximum concerns the signature lists. It is not
an optimality claim for inversion algorithms or for the number of
nonempty cells. At width 32, $F_{33}=3{,}524{,}578$, whereas the word
domain has $2^{32}=4{,}294{,}967{,}296$ members. A signature still incurs
its compatibility and linear-algebra cost.

\section{Reflection targets}
\begin{proposition}\label{prop:reflection}
If $S\in\{2^w-1,2^{w-1}-1\}$, the output list consists only of $k=S$.
Set $\delta=L^{-1}(S)$ and $d=\res{\delta-D}$. If $d$ is even and
$u=d/2$ is supported on $\delta\andb\ell$, then
\begin{equation}\label{eq:reflection}
\mathcal A_{D,S}(L)=\{a:a\andb(\delta\andb\ell)=u\},
\qquad |\mathcal A_{D,S}(L)|=2^{w-\wt(\delta\andb\ell)}.
\end{equation}
Otherwise the fiber is empty.
\end{proposition}
\begin{proof}
The lower $w-1$ target bits are one. Starting from zero carry,
\eqref{eq:outbit} therefore forces $k_i=1$ and keeps the carry zero
through those positions. The top bit then gives $k=S$. The output
mask $\overline{k}\andb\ell$ vanishes, leaving only the input mask.
Its one-bit positions fix exactly $\wt(\delta\andb\ell)$ independent
coordinates.
\end{proof}

A concrete 32-bit example uses the linear map
\begin{equation}\label{eq:map32}
L(x)=\operatorname{ROR}_{17}(x)\oplus\operatorname{ROR}_{19}(x)
       \oplus\operatorname{SHR}_{10}(x).
\end{equation}
Here rotations are cyclic within 32 bits and the right shift fills with
zeros. The supplied verifier reduces its binary matrix to the identity
and obtains
\[
L^{-1}(\mathtt{7fffffff})=\mathtt{59e0cce9}.
\]
Take $D=\mathtt{59e0cce9}$ and $S=\mathtt{7fffffff}$.
Then $u=0$, and the single condition
$a\andb\mathtt{59e0cce9}=0$ describes all $2^{16}=65{,}536$ solutions.
The coefficient and inverse matrices are supplied with the example.

More generally, let $L_w$ keep the low bit and XOR it into every higher
bit. This map is an involution. For $D=1$ and $S=2^w-1$,
$L_w^{-1}(S)=1$, so the complete fiber consists of the even words and
has $2^{w-1}$ members. This family illustrates the difference between
compact representation and explicit enumeration.

\section{A complete eight-bit example}
Represent a binary-linear map by its row masks, with row zero producing
the least significant output bit. Consider
\[
(r_0,\ldots,r_7)=(34,68,134,14,213,186,122,142),\qquad D=30,
\quad S=146.
\]
The row matrix is invertible. There are 13 input signatures and 26 output
signatures. Selecting the input list yields the following five nonempty
cells. Every table entry is a decimal word.
\begin{center}
\begin{tabular}{rrrrl r}
\toprule
$k$&$\delta$&Base&Dimension&Basis&Members\\
\midrule
138&62&209&0&$\varnothing$&1\\
108&254&113&1&129&2\\
18&38&140&1&1&2\\
140&34&15&2&25, 205&4\\
106&226&98&3&8, 17, 132&8\\
\bottomrule
\end{tabular}
\end{center}
Their disjoint union contains 17 points. Both complete signature lists,
the eight discarded input cases, all bases, and all members are
reproducible from \texttt{evidence/figure.json}. The accompanying
three-dimensional figure uses these exact data.

\section{Certificates and verification}
For a nonempty cell, a certificate consists of $k,\delta$, a base, and
independent direction vectors. A checker verifies $L(\delta)=k$, both
mask equations at the base, and the corresponding homogeneous equations
on every direction. It also checks that the basis dimension equals
$w-\rank(A)$, where $A$ is the cell's equation matrix. These conditions
certify the entire cell, including its dimension. Completeness across
cells additionally requires the exact signature traversal proved above.

The publication verifier uses a direct bit-by-bit forward evaluator that
does not use the inverse construction. It compares complete fibers for
every $D,S$ at widths 1 through 7, with the identity and two mixed
invertible matrices at each width. The result is 65,532 fiber comparisons
and 196,596 solver calls across input, output and automatic enumeration.
Every returned set agrees, including empty fibers and multiplicities.

A separate check constructs the feasible signature sets directly from
addition at widths 1 through 8. All 1,020 comparisons agree with the
carry automata. Exhaustive maxima through width 16 attain $F_{w+1}$
on both sides. The two 32-bit reflection examples are expanded and
checked in full. Wide certificates through 512 bits verify the
closed-form family without expanding its output.

The source package contains \texttt{atlas.py}, \texttt{verify.py}, the
exact examples, and the figure builder. The proofs establish the general
statements; finite checks establish the behavior of these implementations.
The reported complexity counts compact affine output. No running-time
or energy claim follows from the signature count alone.

\begin{thebibliography}{9}
\bibitem{mouha}
N. Mouha, V. Velichkov, C. De Canni\`ere, and B. Preneel.
The differential analysis of S-functions.
In \emph{Selected Areas in Cryptography}, LNCS 6544, pp.~36--56.
Springer, 2011. \href{https://doi.org/10.1007/978-3-642-19574-7_3}{doi:10.1007/978-3-642-19574-7\_3}.
\bibitem{velichkov}
V. Velichkov, N. Mouha, C. De Canni\`ere, and B. Preneel.
The additive differential probability of ARX.
In \emph{Fast Software Encryption}, LNCS 6733, pp.~342--358.
Springer, 2011. \href{https://doi.org/10.1007/978-3-642-21702-9_20}{doi:10.1007/978-3-642-21702-9\_20}.
\bibitem{qin}
H. Qin and B. Wu.
Towards non-independence of modular additions in searching differential
trails of ARX ciphers: new automatic methods with application to SPECK
and Chaskey. 2022. \href{https://arxiv.org/abs/2203.09741}{arXiv:2203.09741}.
\end{thebibliography}
\end{document}
